% Original Girard1987 trip/empire proofs; no modern splitting dependency.
\section{Original trip intervals and maximal empires}

\begin{annexitem}{annex:p07.girard-empire.local-blocks}
 {thm:p07-girard-empire-local-blocks}{All six local endpoint matchings}
\begin{proof}
Write \(c,d,r\) for the heads \(C^+,D^+,R^-\), and also for the
corresponding tails \(C^-,D^-,R^+\).
Remove the three successor arcs leaving those tails.
Because the original permutation is one cycle, cutting it at three
distinct places gives three disjoint paths, with these heads and tails,
covering every gate. The removed-tail set and all remaining successors
are identical at \(L\) and \(R\).
Let \(h\) send a head name to the name of the tail at the end of its
path. Let \(g_\tau\) send each tail name to the next head name in
the local routing. The cycles of \(g_\tau h\) are precisely the
orbits obtained by joining these three paths. Thus both compositions
must be three-cycles, since both settings have one long trip.

For tensor, \(g_L=(c\,r\,d)\) and \(g_R=(c\,d\,r)\).
If \(h\) is a transposition, both compositions are odd, hence cannot
be three-cycles. If \(h\) is either three-cycle, one composition is
the identity, giving three orbits. Only \(h=1\) remains, and it does
give a three-cycle in both settings. The paths are therefore
\(C^+\ldots C^-\), \(D^+\ldots D^-\), \(R^-\ldots R^+\).
The routing rows give their stated orders.

For par, \(g_L=(c\,r)\) and \(g_R=(d\,r)\).
If \(h=1\) or either three-cycle, the compositions are odd and fail.
If \(h=(c\,r)\), the \(L\) composition is the identity; if
\(h=(d\,r)\), the \(R\) composition is the identity.
The remaining \(h=(c\,d)\) gives a three-cycle at both settings.
Thus its paths are \(C^+\ldots D^-\), \(D^+\ldots C^-\),
\(R^-\ldots R^+\). Again the routing gives the stated orders.
This accounts for the identity, all three transpositions and both
three-cycles in each link type. It establishes the passage orders
without an empire or sequentialization premise.
\end{proof}
\end{annexitem}

\begin{annexitem}{annex:p07.girard-empire.interval-cases}
 {thm:p07-girard-empire-interval-cases}{Both same-path orientations and all distinct-path positions}
\begin{proof}
If \(a^+,a^-\) lie in one path and \(a^+\) is earlier internally,
their forward interval is an internal segment of that path in both
orders. It cannot contain both endpoints of that path, and contains
neither endpoint pair of either other path.
If \(a^-\) is earlier than \(a^+\), the forward interval goes from
\(a^+\) to the path's end, traverses the other two paths, and returns
from this path's start to \(a^-\). It contains all six endpoints,
regardless of the order of the other paths.

If the gates lie in distinct paths \(X,Y\), with \(a^+\) in \(X\)
and \(a^-\) in \(Y\), the start endpoint of \(X\) and end endpoint
of \(Y\) are excluded in both orders. Thus neither \(X\)'s nor
\(Y\)'s full endpoint pair can be included.
The remaining path \(Z\) is traversed between these interior gates
in exactly one of the two cyclic orders, not both. Therefore its
endpoint pair cannot be included in both. Renaming \(X,Y,Z\)
covers all six ordered choices of two different paths.
Only the same-path, backwards internal order meets the hypothesis,
and in that case the asserted inclusion of all endpoints was proved.
\end{proof}
\end{annexitem}

\begin{annexitem}{annex:p07.girard-empire.closure}
 {thm:p07-girard-empire-closure}{All original facts, with the distinguished and root cases}
\begin{proof}
Fix the tensor \(K\) with premises \(A,B\) and root \(T\).
Its local-block normal form gives disjoint paths
\(I_\sigma(A),I_\sigma(B),[T^-,T^+]\).
The first two are unchanged when only \(K\) is flipped.
Thus \(A\in eA\), the two empires are disjoint, and neither gate
of \(T\) is in either premise interval. This proves (i)--(ii).

For (iii), the axiom routing puts \(C'^-\) immediately after \(C^+\)
and \(C'^+\) immediately before \(C^-\).
If \(C'=A\), the conclusion is (i).
If \(C=A\), these adjacent gates are just inside the two endpoints
of \(I_\sigma(A)\). Otherwise \(C^+,C^-\) are internal points of
that interval, and the adjacent companion gates cannot cross its
endpoints unless \(C'=A\), the case already handled.
Both companion gates are inside at every setting, proving (iii).

For (iv), a tensor root equal to \(T\) cannot satisfy the antecedent
by (ii). A tensor root equal to \(A\) is a separate direct case:
its local root path runs from \(A^-\) to \(A^+\); the complementary
forward interval \(A^+\ldots A^-\) contains both complete premise
paths at both switch settings. Hence both premises are in \(eA\).
For every other tensor link \(L\), \(A\) is neither its root nor
one of its premises: if it were a premise, its unique consumer would
be both \(L\) and \(K\). Therefore \(A^+,A^-\) are non-endpoint
gates of the three paths for \(L\).
Fix every other switch arbitrarily. Since its root belongs to \(eA\),
both root gates belong to the \(A\)-interval at both \(L\) and \(R\).
They are the endpoints of the root path. The interval-cases lemma
puts all six endpoints in that interval in both orders. In particular
both gates of both premises are inside. The other settings were
arbitrary, so this is the universal membership required by (iv).

For (v), a par cannot be the distinguished tensor \(K\).
If its root is \(A\), the same complementary-root-path argument
contains both cross-paired premise paths and hence all four premise
gates. Otherwise \(A\) is not a premise, by unique consumption.
Apply the interval-cases lemma to the root-path endpoints at both
par settings. Its conclusion includes \(C^+,D^-\) and \(D^+,C^-\),
proving both premise memberships at every setting.

For (vi), the link is explicitly different from \(K\).
If its root is \(A\), it is already in \(eA\).
Otherwise \(A\) is neither root nor premise. If \(C\in eA\),
the endpoints of the tensor's \(C\)-path are in the \(A\)-interval
at both settings. The interval-cases lemma includes the root-path
endpoints at both settings. Arbitrary other choices give the claimed
root membership. If \(D\in eA\), use the \(D\)-path instead.
There is no extension to \(K\): its premise \(A\) is in \(eA\)
but its root \(T\) is outside it, as (ii) already proves.

For (vii), a root equal to \(A\) is immediate. Otherwise unique
consumption excludes \(A=C,D\). The two premise memberships put
all four gates \(C^+,C^-,D^+,D^-\) inside at both settings.
In particular both endpoints \(C^+,D^-\) of one par path are inside.
The interval-cases lemma includes both root-path endpoints.
This proves root membership universally. Exchanging the distinguished
premises \(A,B\) gives exactly the symmetric seven facts.
\end{proof}
\end{annexitem}

\begin{annexitem}{annex:p07.girard-empire.realization}
 {thm:p07-girard-empire-realization}{Every possible exit and the arbitrary internal settings}
\begin{proof}
Put \(E=eA\). At a par with exactly one premise in \(E\), prescribe
the position making that premise unselected. Thus its minus gate
returns to its own plus gate: right for an inside left premise,
left for an inside right premise. Choose all other switches arbitrarily.
Follow the trip from \(A^+\) until \(A^-\).
The following list proves that every visited gate belongs to \(\widehat E\).

At a plus gate, its producing axiom sends it to the dual companion's
minus gate, inside by closure (iii). If its producer is tensor or par,
closure (iv) or (v) puts both premises in \(E\), so either possible
departure goes to an inside plus gate.
At a minus gate of a terminal occurrence, the return is to its own
plus gate. At a tensor consumer different from the distinguished
tensor, closure (vi) puts the root in \(E\), and closure (iv) puts
the other premise in \(E\). The tensor routing can go to that root's
minus gate or that other premise's plus gate; both are inside.
At the distinguished tensor the only inside premise is \(A\), since
\(B\in eB\) and \(eA\cap eB=\varnothing\); the traversal stops at
\(A^-\) before taking that exit.

At a par consumer with both premises in \(E\), closure (vii) puts
its root in \(E\); both the selected departure to the root minus
gate and the unselected reflection stay inside.
With exactly one inside premise the prescribed position reflects
that premise to itself. These exhaust the original successor clauses.
Thus the entire interval is contained in \(\widehat E\).
The full-orbit hypothesis guarantees that \(A^-\) is reached.
The reverse inclusion holds by the universal definition of \(eA\).
Hence the interval equals \(\widehat E\).

Only the one-inside-premise par positions were prescribed. Every
other position remained arbitrary throughout the exit proof.
In particular a par producing an inside root has both premises inside
by (v), and is not a boundary par. Its position may be chosen along
an inside producer path, as may every tensor position.
\end{proof}
\end{annexitem}

\begin{annexitem}{annex:p07.girard-empire.ancestor}
 {thm:p07-girard-empire-ancestor}{A freely selected producer path and genuine universal inclusion}
\begin{proof}
Producer closure (iv)--(v), starting from \(A\in eA\), puts every
occurrence along a backwards producer-premise chain in \(eA\).
Let \(R=C\otimes D\) be \(A\) or upstream of it.
Choose the free switches along a producer path from \(A\) through
\(R\) to \(D\) so that the trip begins
\(A^+,\ldots,R^+,D^+\), using only plus gates on that initial path.
A tensor root can select either premise by its two positions; a par
root can likewise select either. No par on this path is a boundary par:
its root is in \(eA\), so both premises are there.
The realization theorem therefore permits these choices while still
giving \(I_\sigma(A)=\widehat{eA}\).

Both gates of \(D\) are in that interval. Its minus gate cannot have
appeared before \(D^+\) on the chosen initial path, which used only
plus gates. Thus the whole forward interval \(I_\sigma(D)\) lies in
\(I_\sigma(A)\), without wrapping past its endpoint.
For any \(F\in eD\), the definition puts both gates of \(F\) in
\(I_\sigma(D)\) at this particular setting. They therefore belong to
\(\widehat{eA}\), implying \(F\in eA\).
This is an inclusion of the universal empire \(eD\), not just of a
chosen switching's occurrence subset. Select the path to \(C\)
instead to obtain \(eC\subseteq eA\).
The argument also covers \(R=A\); its initial producer path has
zero steps before choosing that tensor's child, not a fictitious
extra hereditary premise.
\end{proof}
\end{annexitem}

\begin{annexitem}{annex:p07.girard-empire.joint}
 {thm:p07-girard-empire-joint}{Compatible boundary prescriptions and exact covering}
\begin{proof}
Put \(E=eA,F=eB\); they are disjoint by closure (ii).
Prescribe all boundary par positions for realization of \(E\) and of
\(F\). A conflict could occur only at a par prescribed for both.
Its single inside premise for \(E\) and single inside premise for \(F\)
cannot be the same occurrence, because \(E\cap F=\varnothing\).
Hence such a conflict is exactly a par with one premise in each set.
The assumed absence of crossing pars therefore makes the prescriptions
compatible. Complete all other settings arbitrarily.
The realization theorem applied twice gives
\(I_\sigma(A)=\widehat E\) and \(I_\sigma(B)=\widehat F\)
at this same setting.

The local tensor normal form partitions all gates into these two
premise paths and the root path from \(T^-\) to \(T^+\).
Since \(T\) is terminal, that last successor is a direct return,
so its path consists of exactly the two root gates.
Every occurrence is consequently in \(E\), \(F\), or is \(T\).
The three sets are disjoint by closure (ii), establishing exact covering.
Thus if covering fails, one of the actual par prescriptions conflicts
and supplies an actual crossing par.
\end{proof}
\end{annexitem}

\begin{annexitem}{annex:p07.girard-empire.crossing}
 {thm:p07-girard-empire-crossing}{The two larger-empire settings, without a nesting assumption}
\begin{proof}
First suppose \(C\) is the left premise in \(eA\) and \(D\) the
right premise in \(eB\). The other ordered case is treated below.
The root \(S\) is in neither old empire: if it were in \(eA\),
par producer closure would put \(D\) there, contradicting disjointness;
the symmetric argument excludes \(eB\).

Producer closure from \(J\in eJ\) puts \(S\), and then \(C,D\),
in \(eJ\). Choose free switches along the producer path from \(J\)
to \(S\) so the trip starts with plus gates
\(J^+,\ldots,S^+\), and choose \(L\) at \(S\).
All these producer-root occurrences are inside \(eJ\), so none of
their pars is a reflecting boundary par.
The realization theorem allows these choices while giving
\(I_\sigma(J)=\widehat{eJ}\).

At the selected par the local block form is
\[
 S^+,\,C^+,\,\ldots,D^-,\,D^+,\,\ldots,C^-,\,S^-.
\]
This whole segment is within \(I_\sigma(J)\): \(S^+\) was reached
on the initial plus-only path, so \(S^-\) cannot have occurred earlier;
both gates of \(S\) belong to the realized interval.
This includes the case \(S=J\), where they are its two endpoints.

The old tensor intervals \(I_\sigma(A)\) and \(I_\sigma(B)\) are
disjoint. Both \(C\) gates are in the former and both \(D\) gates
in the latter, by their universal empire memberships.
On a directed circle, a proper interval containing \(C^+,C^-\)
but excluding the intervening \(D^-,D^+\) must run around the
opposite side from \(C^-\) back to \(C^+\).
Its complement is a single interval within the arc \(C^+\ldots C^-\).
The disjoint interval \(I_\sigma(B)\), which contains both \(D\) gates,
lies entirely in that complement. Therefore
\[
 \widehat{eB}\subseteq I_\sigma(B)
       \subseteq [C^+,C^-]_\sigma
       \subseteq I_\sigma(J)=\widehat{eJ}.
\]
The first inclusion is the definition of the universal \(eB\);
the middle inclusions are a statement about this actual chosen trip.
Hence \(eB\subseteq eJ\), without mistaking one interval for an empire.

Now use a second setting realizing \(eJ\), with the same selected
producer path to \(S^+\), but position \(R\) at \(S\).
That free choice gives the local segment
\(S^+,D^+,\ldots,C^-,C^+,\ldots,D^-,S^-\).
The same circle argument with \(C,D\) and the old empires exchanged
puts \(I_\sigma(A)\) inside \(D^+\ldots D^-\), and then inside
\(I_\sigma(J)\). It yields \(eA\subseteq eJ\).
These were two settings, each used to bound a universal empire.

If the ordered premises instead have left \(C\in eB\) and right
\(D\in eA\), apply those same two circle arguments with the old
empire names exchanged, using the actual \(L\) and \(R\) ports as
above. No connective's ordered premise positions are changed.
Both old empires are again contained in \(eJ\).
The inclusion of their union is strict because \(S\in eJ\) and
\(S\) belongs to neither old empire.
Finally \(eJ\subseteq eH\cup eJ\), the union attached to the
terminal tensor \(U\); this is valid whether \(J\) is its left
or its right premise.
\end{proof}
\end{annexitem}

\begin{annexitem}{annex:p07.girard-empire.splitting}
 {thm:p07-girard-empire-splitting}{Terminality and covering are consequences, not premises}
\begin{proof}
There are finitely many tensor links and at least one by hypothesis.
Their empire unions therefore have a maximal member under inclusion.
Choose its link \(T=A\otimes B\).

If \(T\) is not terminal, follow its unique consumer and then successive
consumers to a terminal occurrence. Formula complexity strictly
increases at every such step, so finiteness guarantees termination.
The final producer is a connective, and the hypothesis excludes a
terminal par; hence it is a terminal tensor \(U\).
Let \(J\) be the premise of \(U\) equal to or downstream of \(T\).
The ancestor corollary applied with distinguished empire \(eJ\)
gives \(eA\cup eB\subseteq eJ\).
Producer closure puts \(T\in eJ\), whereas the original local tensor
blocks put \(T\notin eA\cup eB\).
Consequently the inclusion is strict, and \(eJ\) is contained in the
union of the two empires of \(U\), contradicting maximality.
Thus \(T\) is terminal.

If its empire union and root do not cover the structure, the joint
realization lemma produces a crossing par \(S\).
Its root is in neither empire by par producer closure.
It cannot be terminal, by hypothesis. Following its consumers as
above reaches a terminal tensor \(U\), with premise \(J\) equal to
or downstream of \(S\). The proved crossing lemma gives
\(eA\cup eB\subsetneq eJ\), contained in the union at \(U\).
This again contradicts maximality. Therefore covering holds.
Disjointness and the root's exclusion were proved in closure (ii).
No terminal splitter or sequentialization was assumed in selecting
or enlarging the maximal empire union.
\end{proof}
\end{annexitem}

\begin{annexitem}{annex:p07.girard-empire.restriction}
 {thm:p07-girard-empire-restriction}{Every link and every restricted setting}
\begin{proof}
Put \(E=eA,F=eB\), so the splitting theorem gives the disjoint
partition \(E,F,\{T\}\).
If an occurrence in \(E\) has an axiom producer, (iii) includes the
other axiom output. If it has a tensor or par producer, (iv)--(v)
include both premises. Thus every retained occurrence has its entire
producer link in the restriction.
A tensor consuming a retained premise is either the removed \(T\),
consuming \(A\), or a different tensor. In the latter case (vi),
followed by (iv), puts its root and other premise in \(E\).
A par cannot consume one premise in \(E\) and one outside:
the other premise is then in \(F\), since terminal \(T\) is not
a premise of any link. Its root would belong to neither \(E\) nor
\(F\) by producer closure and disjointness, and cannot equal tensor
\(T\). This contradicts the partition. A par with both premises
in \(E\) has its root there by (vii).
Hence all retained consumers are complete, except the deleted \(T\).
The restriction is a valid original proof structure, and \(A\) is
newly terminal. It is nonempty because \(A\in E\).

Give its retained switches any setting \(\tau\), and extend that
setting arbitrarily to the other original links.
There are no boundary pars with one premise in \(E\), as just proved.
Thus the realization theorem's reflecting prescriptions impose no
condition on the retained settings, and its complete exit argument
gives \(I_\sigma(A)=\widehat E\).
All successors on this interval agree with those of the restriction
until its last gate \(A^-\). The restriction replaces only the old
exit through \(T\) by the new terminal return \(A^-\mapsto A^+\).
Its orbit consequently traverses every retained gate exactly once.
As \(\tau\) was arbitrary, the restriction is trip-correct at every
setting. Exchange \(E,F\) and \(A,B\) for the other restriction.
This supplies the correctness step omitted in the printed continuation
of the original sequentialization proof.
\end{proof}
\end{annexitem}

\begin{annexitem}{annex:p07.girard-empire.sequentialization}
 {thm:p07-girard-empire-sequentialization}{All original induction cases and graph reconstruction}
\begin{proof}
Induct on the number of links.
With exactly one link, it is an axiom: a tensor or par would require
premise occurrences produced by another link, since formula complexity
excludes self-production. Take the general-formula sequent axiom.

If a terminal par has premises \(C,D\) and root \(S\), remove that
link and root occurrence, exposing \(C,D\) as terminal.
The result is a nonempty original proof structure with one fewer link.
For any setting of it, extend to the original by position \(L\) at
the removed par. Since \(S\) was terminal, its two consecutive ports
occur in the full orbit as
\(C^-\mapsto S^-\mapsto S^+\mapsto C^+\).
Removing the two \(S\) ports replaces precisely this chain by
\(C^-\mapsto C^+\). The \(D\) return is unchanged.
Thus the entire reduced orbit is still one long trip for every
setting of the reduced structure. Apply induction and then the par
sequent rule to the distinguished conclusions \(C,D\).
The occurrence construction restores exactly that removed link.

Otherwise there is no terminal par. If more than one link exists,
there must be a connective link: two or more axiom-only components
have separate trip orbits, contrary to correctness.
Follow a connective's consumer chain to a terminal connective;
finiteness and increasing formula complexity guarantee one.
It is a tensor. The maximal-empire splitting theorem and the
restriction lemma give two smaller nonempty correct structures,
with distinguished conclusions \(A,B\), whose links partition all
remaining links. Inductively obtain their Mix-free proofs and apply
tensor to these disjoint premise proofs. Their other terminal
occurrences are retained, and the tensor restores root \(T\).
This reconstructs exactly the original incidence structure.
The rules used are general-formula axiom, exchange, par and tensor;
no Mix, modern splitter or alternative correctness criterion enters.
\end{proof}
\end{annexitem}

